% Options for packages loaded elsewhere
% Options for packages loaded elsewhere
\PassOptionsToPackage{unicode}{hyperref}
\PassOptionsToPackage{hyphens}{url}
\PassOptionsToPackage{dvipsnames,svgnames,x11names}{xcolor}
%
\documentclass[
  11pt,
]{article}
\usepackage{xcolor}
\usepackage[margin=1in]{geometry}
\usepackage{amsmath,amssymb}
\setcounter{secnumdepth}{-\maxdimen} % remove section numbering
\usepackage{iftex}
\ifPDFTeX
  \usepackage[T1]{fontenc}
  \usepackage[utf8]{inputenc}
  \usepackage{textcomp} % provide euro and other symbols
\else % if luatex or xetex
  \usepackage{unicode-math} % this also loads fontspec
  \defaultfontfeatures{Scale=MatchLowercase}
  \defaultfontfeatures[\rmfamily]{Ligatures=TeX,Scale=1}
\fi
\usepackage{lmodern}
\ifPDFTeX\else
  % xetex/luatex font selection
\fi
% Use upquote if available, for straight quotes in verbatim environments
\IfFileExists{upquote.sty}{\usepackage{upquote}}{}
\IfFileExists{microtype.sty}{% use microtype if available
  \usepackage[]{microtype}
  \UseMicrotypeSet[protrusion]{basicmath} % disable protrusion for tt fonts
}{}
\makeatletter
\@ifundefined{KOMAClassName}{% if non-KOMA class
  \IfFileExists{parskip.sty}{%
    \usepackage{parskip}
  }{% else
    \setlength{\parindent}{0pt}
    \setlength{\parskip}{6pt plus 2pt minus 1pt}}
}{% if KOMA class
  \KOMAoptions{parskip=half}}
\makeatother
% Make \paragraph and \subparagraph free-standing
\makeatletter
\ifx\paragraph\undefined\else
  \let\oldparagraph\paragraph
  \renewcommand{\paragraph}{
    \@ifstar
      \xxxParagraphStar
      \xxxParagraphNoStar
  }
  \newcommand{\xxxParagraphStar}[1]{\oldparagraph*{#1}\mbox{}}
  \newcommand{\xxxParagraphNoStar}[1]{\oldparagraph{#1}\mbox{}}
\fi
\ifx\subparagraph\undefined\else
  \let\oldsubparagraph\subparagraph
  \renewcommand{\subparagraph}{
    \@ifstar
      \xxxSubParagraphStar
      \xxxSubParagraphNoStar
  }
  \newcommand{\xxxSubParagraphStar}[1]{\oldsubparagraph*{#1}\mbox{}}
  \newcommand{\xxxSubParagraphNoStar}[1]{\oldsubparagraph{#1}\mbox{}}
\fi
\makeatother


\usepackage{longtable,booktabs,array}
\usepackage{calc} % for calculating minipage widths
% Correct order of tables after \paragraph or \subparagraph
\usepackage{etoolbox}
\makeatletter
\patchcmd\longtable{\par}{\if@noskipsec\mbox{}\fi\par}{}{}
\makeatother
% Allow footnotes in longtable head/foot
\IfFileExists{footnotehyper.sty}{\usepackage{footnotehyper}}{\usepackage{footnote}}
\makesavenoteenv{longtable}
\usepackage{graphicx}
\makeatletter
\newsavebox\pandoc@box
\newcommand*\pandocbounded[1]{% scales image to fit in text height/width
  \sbox\pandoc@box{#1}%
  \Gscale@div\@tempa{\textheight}{\dimexpr\ht\pandoc@box+\dp\pandoc@box\relax}%
  \Gscale@div\@tempb{\linewidth}{\wd\pandoc@box}%
  \ifdim\@tempb\p@<\@tempa\p@\let\@tempa\@tempb\fi% select the smaller of both
  \ifdim\@tempa\p@<\p@\scalebox{\@tempa}{\usebox\pandoc@box}%
  \else\usebox{\pandoc@box}%
  \fi%
}
% Set default figure placement to htbp
\def\fps@figure{htbp}
\makeatother





\setlength{\emergencystretch}{3em} % prevent overfull lines

\providecommand{\tightlist}{%
  \setlength{\itemsep}{0pt}\setlength{\parskip}{0pt}}



 


\usepackage{amsmath,amssymb,enumitem,booktabs}
\setlist{itemsep=2pt,topsep=3pt}
\usepackage{fancyhdr}\pagestyle{fancy}\fancyhf{}
\fancyhead[L]{\small DS701 Fall 2026}\fancyhead[R]{\small Homework 3}
\fancyfoot[C]{\small\thepage}
\renewcommand{\headrulewidth}{0.4pt}
\usepackage[breakable]{tcolorbox}
% --- answer-space macros (problems PDF only; no-ops in the solutions PDF) ---
% \answerspace{2in}  vertical room for handwritten work, sized from the solution
% \blank             a short fill-in rule; \blank[1.5in] for a wider one
% \answerpage        start the next problem on a fresh page (Gradescope page matching)
\newcommand{\answerspace}[1]{\par\vspace*{#1}}  % \vspace* so room is not dropped at a page break
\newcommand{\blank}[1][1.2in]{\rule{#1}{0.4pt}}
\newcommand{\answerpage}{\clearpage}
\newtcolorbox{solutionbox}{breakable,colback=blue!4,colframe=blue!45!black,title=Solution,fonttitle=\bfseries,left=4pt,right=4pt,top=3pt,bottom=3pt}
\makeatletter
\@ifpackageloaded{caption}{}{\usepackage{caption}}
\AtBeginDocument{%
\ifdefined\contentsname
  \renewcommand*\contentsname{Table of contents}
\else
  \newcommand\contentsname{Table of contents}
\fi
\ifdefined\listfigurename
  \renewcommand*\listfigurename{List of Figures}
\else
  \newcommand\listfigurename{List of Figures}
\fi
\ifdefined\listtablename
  \renewcommand*\listtablename{List of Tables}
\else
  \newcommand\listtablename{List of Tables}
\fi
\ifdefined\figurename
  \renewcommand*\figurename{Figure}
\else
  \newcommand\figurename{Figure}
\fi
\ifdefined\tablename
  \renewcommand*\tablename{Table}
\else
  \newcommand\tablename{Table}
\fi
}
\@ifpackageloaded{float}{}{\usepackage{float}}
\floatstyle{ruled}
\@ifundefined{c@chapter}{\newfloat{codelisting}{h}{lop}}{\newfloat{codelisting}{h}{lop}[chapter]}
\floatname{codelisting}{Listing}
\newcommand*\listoflistings{\listof{codelisting}{List of Listings}}
\makeatother
\makeatletter
\makeatother
\makeatletter
\@ifpackageloaded{caption}{}{\usepackage{caption}}
\@ifpackageloaded{subcaption}{}{\usepackage{subcaption}}
\makeatother
\usepackage{bookmark}
\IfFileExists{xurl.sty}{\usepackage{xurl}}{} % add URL line breaks if available
\urlstyle{same}
\hypersetup{
  pdftitle={DS701 Homework 3: k-means and Linkage},
  colorlinks=true,
  linkcolor={blue},
  filecolor={Maroon},
  citecolor={Blue},
  urlcolor={Blue},
  pdfcreator={LaTeX via pandoc}}


\title{DS701 Homework 3: k-means and Linkage}
\usepackage{etoolbox}
\makeatletter
\providecommand{\subtitle}[1]{% add subtitle to \maketitle
  \apptocmd{\@title}{\par {\large #1 \par}}{}{}
}
\makeatother
\subtitle{Released Mon Sep 21 · due Mon Sep 28, 11:59 pm on Gradescope}
\author{}
\date{}
\begin{document}
\maketitle


\noindent\textbf{Name:}\ \rule{2.0in}{0.4pt}\qquad\textbf{BU ID:}\ \rule{1.2in}{0.4pt}
\vspace{6pt}

\textbf{How this works.} Answer \textbf{by hand on paper} (or a tablet
with a stylus), show your work, then scan and upload to Gradescope,
matching each problem to its pages. This is deliberate practice for the
closed-notes quizzes, which are in the same style --- so do it without a
computer or an AI assistant. It is \textbf{participation-graded}: a
serious attempt at every problem earns full credit, and worked solutions
are released the day after the deadline. Budget 1--2 hours.

\textbf{Covers:} Lecture 03 (Clustering I: \(k\)-means) and Lecture 04
(Clustering II: hierarchical).

\subsection{Problem 1: Lloyd's algorithm by
hand}\label{problem-1-lloyds-algorithm-by-hand}

\emph{Practices: running \(k\)-means step by step, tracking the
objective, and seeing a bad local minimum.}

Six one-dimensional data points: \[x = 0,\ 5,\ 6,\ 7,\ 12,\ 15 .\]
Cluster them with \(k\)-means, \(K = 2\), using Euclidean distance (on a
line, \(\|x-c\|^2 = (x-c)^2\)). The objective is the WCSS / inertia,
\(\sum_k \sum_{x\in C_k}(x-c_k)^2\).

\begin{enumerate}
\def\labelenumi{(\alph{enumi})}
\tightlist
\item
  Start from the centers \(c_1 = 12\), \(c_2 = 15\). Run
  \textbf{iteration 1}: assign every point to its nearest center, then
  update each center to the mean of its cluster. Report the two
  clusters, the WCSS \emph{right after the assignment step} (using the
  old centers), the new centers, and the WCSS \emph{right after the
  update step} (using the new centers, same assignment).
\end{enumerate}

\begin{longtable}[]{@{}lll@{}}
\toprule\noalign{}
Iteration 1 & Cluster 1 & Cluster 2 \\
\midrule\noalign{}
\endhead
\bottomrule\noalign{}
\endlastfoot
Points assigned & \blank[1.6in] & \blank[1.6in] \\
WCSS after assignment (old centers) & \blank[1.6in] & \\
New centers & \blank[1.6in] & \blank[1.6in] \\
WCSS after update (new centers) & \blank[1.6in] & \\
\end{longtable}

\answerspace{2.5in}

\begin{enumerate}
\def\labelenumi{(\alph{enumi})}
\setcounter{enumi}{1}
\tightlist
\item
  Run \textbf{iteration 2} the same way and report the same four things.
  Then state whether a third assignment step would change anything, and
  list the WCSS values you have seen, in order. Did the objective ever
  go up?
\end{enumerate}

\begin{longtable}[]{@{}lll@{}}
\toprule\noalign{}
Iteration 2 & Cluster 1 & Cluster 2 \\
\midrule\noalign{}
\endhead
\bottomrule\noalign{}
\endlastfoot
Points assigned & \blank[1.6in] & \blank[1.6in] \\
WCSS after assignment (old centers) & \blank[1.6in] & \\
New centers & \blank[1.6in] & \blank[1.6in] \\
WCSS after update (new centers) & \blank[1.6in] & \\
\end{longtable}

Would a third assignment step change anything? \blank[0.8in]\\
WCSS values seen, in order: \blank[2.5in]\quad Ever went up?
\blank[0.8in]

\answerspace{2.5in}

\begin{enumerate}
\def\labelenumi{(\alph{enumi})}
\setcounter{enumi}{2}
\tightlist
\item
  Now start instead from \(c_1 = 0\), \(c_2 = 5\). Run Lloyd's algorithm
  until the centers stop changing (it is short). Report the final
  clusters, centers, and WCSS. Compare with (b): which run found the
  better clustering, and is the other run ``wrong'' from the algorithm's
  point of view?
\end{enumerate}

Final clusters: \blank[2.2in] Centers: \blank[1.2in] WCSS: \blank[0.8in]

\answerspace{3in}

\begin{enumerate}
\def\labelenumi{(\alph{enumi})}
\setcounter{enumi}{3}
\tightlist
\item
  Two sentences: what does (c) imply for how \(k\)-means is run in
  practice, and what would the \(k\)-means++ idea from the lecture have
  done with this data? (Check your claim: run one iteration from the two
  centers \(k\)-means++ would most plausibly pick.)
\end{enumerate}

\answerspace{1.8in}

\answerpage

\subsection{Problem 2: Why k-means stops, and why more clusters always
``help''}\label{problem-2-why-k-means-stops-and-why-more-clusters-always-help}

\emph{Practices: the two one-line facts behind Lloyd's algorithm, and
reading the WCSS correctly when choosing \(K\).}

Write \(J(C, c) = \sum_{k=1}^K \sum_{x\in C_k}\|x-c_k\|^2\) for the
objective, where \(C = (C_1,\dots,C_K)\) is the assignment of points to
clusters and \(c = (c_1,\dots,c_K)\) the centers.

\begin{enumerate}
\def\labelenumi{(\alph{enumi})}
\tightlist
\item
  \textbf{Assignment step.} Hold the centers \(c\) fixed and let \(C'\)
  be the assignment that sends every point to its nearest center. Show
  that \(J(C', c) \le J(C, c)\) for \emph{any} assignment \(C\). (One or
  two lines: compare the two sums term by term.)
\end{enumerate}

\answerspace{1.5in}

\begin{enumerate}
\def\labelenumi{(\alph{enumi})}
\setcounter{enumi}{1}
\tightlist
\item
  \textbf{Update step.} Hold an assignment fixed and look at one cluster
  with points \(x_1,\dots,x_n\) (one-dimensional, to keep the algebra
  short --- the same argument works coordinate by coordinate). Show that
  \(f(c) = \sum_{i=1}^n (x_i - c)^2\) is minimized at
  \(c = \bar x = \frac1n\sum_i x_i\). Either differentiate \(f\) and set
  the derivative to zero (and say why that critical point is a minimum),
  or expand \((x_i - c)^2 = \big((x_i-\bar x) + (\bar x - c)\big)^2\)
  and simplify.
\end{enumerate}

\answerspace{2.5in}

\begin{enumerate}
\def\labelenumi{(\alph{enumi})}
\setcounter{enumi}{2}
\tightlist
\item
  \textbf{Termination.} Using (a), (b), and one further fact --- that
  there are only finitely many ways to assign \(N\) points to \(K\)
  clusters --- explain in a short paragraph why Lloyd's algorithm must
  stop after finitely many iterations (assume, as \texttt{sklearn} does,
  that a point stays where it is when it is equally close to two
  centers). Then say in one sentence what the argument does \emph{not}
  promise about the clustering it stops at.
\end{enumerate}

\answerspace{2in}

\begin{enumerate}
\def\labelenumi{(\alph{enumi})}
\setcounter{enumi}{3}
\tightlist
\item
  \textbf{Choosing \(K\).} Let \(J^*_K\) be the smallest WCSS achievable
  with \(K\) clusters (the global optimum, not a particular run).
  Explain why \(J^*_{K+1} \le J^*_K\) for every \(K\), and what
  \(J^*_N\) is when \(N\) is the number of points. Conclude why ``pick
  the \(K\) with the smallest WCSS'' is not a usable rule, and name the
  criterion the lecture uses instead --- describe in one sentence what
  it looks for.
\end{enumerate}

\(J^*_N =\) \blank[0.8in] \qquad Criterion the lecture uses:
\blank[1.5in]

\answerspace{2.2in}

\answerpage

\subsection{Problem 3: Single vs.~complete linkage on five
points}\label{problem-3-single-vs.-complete-linkage-on-five-points}

\emph{Practices: agglomerative clustering step by step, dendrogram
heights, chaining, and matching a linkage to a data shape.}

Five objects A--E with pairwise distances (a metric):

\begin{longtable}[]{@{}cccccc@{}}
\toprule\noalign{}
& A & B & C & D & E \\
\midrule\noalign{}
\endhead
\bottomrule\noalign{}
\endlastfoot
\textbf{A} & 0 & 2 & 5 & 9 & 12 \\
\textbf{B} & 2 & 0 & 3 & 7 & 10 \\
\textbf{C} & 5 & 3 & 0 & 4 & 8 \\
\textbf{D} & 9 & 7 & 4 & 0 & 5 \\
\textbf{E} & 12 & 10 & 8 & 5 & 0 \\
\end{longtable}

\begin{enumerate}
\def\labelenumi{(\alph{enumi})}
\tightlist
\item
  Perform agglomerative clustering with \textbf{single linkage}. At each
  of the four steps write the current cluster-to-cluster distances that
  matter, the pair merged, and the merge height. Then write the result
  as the \(4\times4\) linkage matrix, rows
  \texttt{{[}idx1,\ idx2,\ dist,\ count{]}}, with A--E as indices
  \(0\)--\(4\) and new clusters numbered \(5, 6, \dots\) in creation
  order.
\end{enumerate}

\begin{longtable}[]{@{}llll@{}}
\toprule\noalign{}
Single linkage & Pair merged & Height & Linkage-matrix row \\
\midrule\noalign{}
\endhead
\bottomrule\noalign{}
\endlastfoot
Step 1 & \blank[1.2in] & \blank[0.7in] & \blank[2in] \\
Step 2 & \blank[1.2in] & \blank[0.7in] & \blank[2in] \\
Step 3 & \blank[1.2in] & \blank[0.7in] & \blank[2in] \\
Step 4 & \blank[1.2in] & \blank[0.7in] & \blank[2in] \\
\end{longtable}

\answerspace{4in}

\begin{enumerate}
\def\labelenumi{(\alph{enumi})}
\setcounter{enumi}{1}
\tightlist
\item
  Repeat with \textbf{complete linkage}.
\end{enumerate}

\begin{longtable}[]{@{}llll@{}}
\toprule\noalign{}
Complete linkage & Pair merged & Height & Linkage-matrix row \\
\midrule\noalign{}
\endhead
\bottomrule\noalign{}
\endlastfoot
Step 1 & \blank[1.2in] & \blank[0.7in] & \blank[2in] \\
Step 2 & \blank[1.2in] & \blank[0.7in] & \blank[2in] \\
Step 3 & \blank[1.2in] & \blank[0.7in] & \blank[2in] \\
Step 4 & \blank[1.2in] & \blank[0.7in] & \blank[2in] \\
\end{longtable}

\answerspace{3.2in}

\begin{enumerate}
\def\labelenumi{(\alph{enumi})}
\setcounter{enumi}{2}
\tightlist
\item
  Sketch both dendrograms (leaves on the x-axis, merge heights on the
  y-axis). For each, say which two clusters you get by cutting the tree
  into two flat clusters. Then explain, in terms of the
  \emph{definitions} of the two linkages, why single linkage swallowed
  A, B, C, D one at a time while complete linkage kept two compact
  groups --- the lecture's word for the single-linkage behavior is
  \emph{chaining}.
\end{enumerate}

\answerspace{3.4in}

Two-cluster cut, single: \blank[2in] \qquad complete: \blank[2in]

\answerspace{1.4in}

\begin{enumerate}
\def\labelenumi{(\alph{enumi})}
\setcounter{enumi}{3}
\tightlist
\item
  For each situation below, name the linkage the lecture would recommend
  and give one sentence of justification: (i) two crescent-shaped
  (``half-moon'') clusters that curve around each other; (ii) several
  roughly round groups of similar size, plus a handful of scattered
  noise points between them.
\end{enumerate}

Situation (i) --- linkage: \blank[1.5in]

\answerspace{1in}

Situation (ii) --- linkage: \blank[1.5in]

\answerspace{1in}




\end{document}
